% #############################################################################
% Abstract Text
% !TEX root = ../main.tex
% #############################################################################
% use \noindent in first paragraph


\noindent Solar powered vehicles, particularly boats, need effective energy conversion in rapidly changing conditions. The \ac{mppt} charger plays a critical role in energy conversion, ensuring the panels operate at their \ac{mpp}. This work focuses on the design, simulation, and testing of a fast response \ac{mppt} charger for \ac{tsb}.

The developed prototype was designed to charge a lithium-ion battery of 43.2~V nominal voltage. The prototype is based on a synchronous boost converter operating at 100~kHz controlled by a \ac{mcu}. 

Different \ac{mppt} techniques were analyzed to identify a suitable solution for the dynamic operating conditions of the vessel. Balancing tracking speed, accuracy and complexity led to the implementation of a \ac{peo} algorithm. The converter and the \ac{pv} model were simulated in LTspice, and the tracker was tuned in MATLAB/Simulink. A custom \ac{pcb} was designed, manufactured and soldered by hand. Firmware on the \ac{mcu} implements \ac{peo} with \acs{cc}-\acs{cv} charging, \ac{can} telemetry and a computer application for configuration.

Laboratory tests gave a peak efficiency of 96.84\%, with tracking typically below 100~ms and 200~ms in the most demanding case, remaining below the 0.5~s and 1~s targets. Conducted noise and a 3~h thermal run stayed within safe limits. The charger followed the \acs{cc}-\acs{cv} profile in the laboratory and in an on-shore field test. At a cost of about 125 €, the board is a dedicated, configurable and low-cost \ac{mppt} charger for \ac{tsb}.

\vspace{0.5cm}